%% kapa: exam=1 round=37 year=2026 period=1 form=A subject=economics subject_ko=경제학원론 seq=9 number=49
%% source: Q-Net 감정평가사 기출문제 (cst00309) artlSeq=5258924 file=raw/1008/5258924/2026년 제37회 감정평가사 1차 1교시.pdf
%% conversion: kordoc 4.24.51 -> tools/gen1.py (auto); review=proofread-sonnet; flags=none
\begin{question}{49}
생산물 시장의 독점기업에 관한 설명으로 옳지 않은 것은? (단, 생산물에 대한 수요곡선은 우하향하는 직선이며, 독점기업은 이윤극대화를 추구한다.)
\begin{choices}
  \item 총수입이 최대인 지점에서 한계수입은 0이다.
  \item 한계수입 > 0이면, 수량 증가시 총수입은 증가한다.
  \item 수요의 가격탄력성이 0.5일 때, 제품가격을 인하하면 총수입은 증가한다.
  \item 수요곡선 상의 한계수입 = 0인 점에서 수요의 가격탄력성은 1이다.
  \item 한계수입 = 한계비용에서 수량을 결정하면 독점이윤이 극대화된다.
\end{choices}
\end{question}
\answer{3}
